\cleardoublepage % memastikan bab baru dimulai di halaman ganjil (kanan)
\chapter{PENDAHULUAN}
\label{chap:pendahuluan}

Bab ini membahas latar belakang pengembangan sistem pelacakan peserta ITB Ultra-Marathon beserta permasalahan yang melatarbelakanginya. Selain itu, bab ini juga memaparkan rumusan masalah, tujuan, batasan masalah, metodologi yang digunakan dalam proses pengembangan sistem, serta sistematika penulisan laporan tugas akhir secara keseluruhan.
% --- Latar Belakang ---
\section{Latar Belakang}
Dalam beberapa tahun terakhir, perkembangan \textit{ultra-marathon} menunjukkan pergeseran dari aktivitas olahraga ekstrem menjadi fenomena gaya hidup global yang semakin populer. Sebagai cabang olahraga lari yang menempuh jarak lebih dari 42,195 kilometer, \textit{ultra-marathon} memerlukan ketahanan fisik dan mental yang luar biasa dari para pesertanya \cite{Ronto2024}. Di Indonesia, salah satu bukti meningkatnya tren ini adalah diselenggarakannya kegiatan tahunan ITB Ultra-Marathon. Sejak pertama kali diselenggarakan pada tahun 2017, ITB Ultra-Marathon terus mengalami peningkatan jumlah peserta dari tahun ke tahun \cite{Hafizh2025}. Hal ini menunjukkan tingginya minat masyarakat terhadap olahraga \textit{ultra-marathon}. Seiring dengan meningkatnya partisipasi pada cabang olahraga ini, kebutuhan akan sistem pelacakan pelari menjadi semakin krusial dalam penyelenggaraan \textit{ultra-marathon}. Sistem pelacakan tidak hanya berperan dalam menjaga keselamatan peserta, tetapi juga dalam memantau performa serta meningkatkan pengalaman pengguna melalui penyajian informasi yang interaktif \cite{Hochreiter2024}.

Dalam konteks ITB Ultra-Marathon, sistem pelacakan dibutuhkan untuk memantau posisi pelari secara individu maupun kelompok. Melalui sistem ini, panitia, tim pendukung, maupun penonton dapat mengetahui lokasi pelari secara \textit{real-time}, termasuk informasi seperti jarak tempuh, waktu, serta elevasi lintasan. Data tersebut sangat penting untuk mendukung proses penjemputan dan pengantaran peserta \textit{ultra-marathon} di titik-titik tertentu sepanjang rute perlombaan. Selain itu, sistem ini juga berperan dalam membantu pelari selama perlombaan berlangsung, misalnya dengan menyediakan informasi rute yang harus diikuti sehingga meminimalkan risiko tersesat. Bagi kategori \textit{relay}, sistem pelacakan turut mempermudah koordinasi pergantian pelari. Dengan demikian, keberadaan sistem pelacakan tidak hanya meningkatkan koordinasi, tetapi juga berkontribusi pada keamanan dan kenyamanan peserta selama \textit{ultra-marathon} berlangsung.

Namun, berdasarkan observasi pada penyelenggaraan ITB Ultra-Marathon di tahun-tahun sebelumnya, sistem pelacakan yang digunakan masih memiliki keterbatasan yang signifikan. Teknologi yang selama ini diandalkan hanya mampu merekam data secara pasif saat pelari melewati titik tertentu. Akibatnya, keberadaan pelari di antara titik tersebut tidak dapat terpantau secara \textit{real-time}. Kesenjangan informasi ini berpotensi menghambat respons terhadap situasi darurat, terutama pada area rute yang terpencil maupun pada saat kondisi perlombaan berlangsung di malam hari.

Selain itu, upaya pelacakan yang pernah dilakukan peserta hanya memanfaatkan aplikasi lari konvensional yang sudah tersedia di pasar. Solusi ini ditemukan tidak efektif untuk kebutuhan manajemen perlombaan berskala besar. Aplikasi lari umum dirancang untuk penggunaan personal dan tidak memiliki fitur integrasi data yang memungkinkan panitia untuk memantau ribuan pelari secara serentak dalam satu antarmuka yang terpusat \cite{higgins2016smartphone}. Penggunaan aplikasi yang terpisah-pisah menyulitkan koordinasi logistik dan sinkronisasi data bagi panitia dan tim pendukung di lapangan. Kesenjangan antara kebutuhan manajemen acara dengan keterbatasan fitur pada aplikasi lari konvensional tersebut menunjukkan perlunya sebuah sistem pelacakan mandiri yang dirancang khusus untuk ekosistem ITB Ultra-Marathon.

Sebagai solusi atas permasalahan tersebut, diusulkan pengembangan sistem pelacak khusus untuk kegiatan ITB Ultra-Marathon yang menitikberatkan pada aspek visualisasi manajemen perlombaan. Penelitian ini difokuskan pada pengembangan sisi \textit{front-end} untuk mentransformasikan data koordinat yang kompleks menjadi informasi visual yang informatif, intuitif, dan responsif. Dengan rancangan antarmuka yang efektif, diharapkan panitia, tim support, penonton, maupun peserta dapat mengakses informasi posisi secara akurat guna menjamin keamanan dan kelancaran operasional selama ITB Ultra-Marathon berlangsung \cite{Caselli2018}.

% --- Rumusan Masalah ---
\section{Rumusan Masalah}
Saat ini belum tersedia sistem yang secara khusus digunakan untuk melakukan pelacakan peserta pada kegiatan ITB Ultra-Marathon. Berdasarkan latar belakang yang telah diuraikan sebelumnya, maka rumusan masalah dalam tugas akhir ini adalah: “Bagaimana mengembangkan sistem \textit{front-end} aplikasi ITB Ultra-Marathon yang mampu menyajikan visualisasi data pelacakan peserta secara \textit{real-time}, responsif, dan mudah digunakan?”

Aplikasi yang dirancang diharapkan mampu menjadi solusi yang memfasilitasi proses pelacakan peserta secara efisien dan \textit{real-time}, sehingga dapat membantu berlangsungnya kegiatan ITB Ultra-Marathon.
% --- Tujuan ---
\section{Tujuan}

Tujuan dari tugas akhir ini adalah mengembangkan dan mengimplementasikansistem \textit{front-end} aplikasi ITB Ultra-Marathon yang mampu menyajikan visualisasi data pelacakan peserta secara \textit{real-time} serta mengintegrasikan sistem \textit{front-end} dengan sistem \textit{back-end} agar data pelacakan dapat ditampilkan secara optimal.
Keberhasilan dari pelaksanaan tugas akhir ini diukur berdasarkan kriteria-kriteria berikut:

\begin{enumerate}
    \item Antarmuka yang dikembangkan mampu memenuhi kebutuhan informasi dan alur penggunaan bagi panitia, tim support, penonton, dan peserta sesuai dengan hasil analisis kebutuhan sistem.
    \item Pengguna dapat memahami alur navigasi serta membaca informasi pelacakan peserta maraton dengan mudah.
    \item Informasi pelacakan peserta maraton dapat visualisasikan secara jelas, akurat, mudah dipahami, serta bersifat responsif terhadap pembaruan data.
    \item Tampilan antarmuka berfungsi dengan baik pada perangkat serta menjaga konsistensi desain.
    \item Terjalinnya integrasi antara \textit{front-end} dan \textit{back-end} yang efektif, di mana data dapat ditampilkan tanpa gangguan signifikan sehingga proses pelacakan dapat dilakukan secara lancar.
\end{enumerate}
% --- Batasan Masalah ---
\section{Batasan Masalah}
Agar pembahasan lebih terarah dan fokus, penelitian ini dibatasi oleh beberapa ruang lingkup sebagai berikut:

\begin{enumerate}
    \item Tugas akhir ini dikerjakan secara berkelompok yang terdiri dari dua orang mahasiswa, yaitu Dinda Thalia Fahira (NIM 18222055) dan Justin Lawrance (NIM 18222006). Adapun pembagian fokus dari kedua mahasiswa tersebut adalah Thalia fokus pada bagian \textit{front-end} aplikasi dan Justin fokus pada bagian \textit{back-end} aplikasi.
    \item Pengembangan aplikasi \textit{tracker} ini dibatasi khusus untuk kebutuhan kegiatan ITB Ultra-Marathon dan tidak mencakup implementasi di luar lingkungan ITB atau skala kegiatan \textit{ultra-marathon} lainnya.
    \item Aplikasi pelacakan ini dikembangkan khusus untuk perangkat \textit{mobile} dengan sistem operasi Android dan tidak mencakup pengembangan untuk sistem operasi iOS.
\end{enumerate}

Batasan masalah ini ditetapkan untuk memastikan bahwa proses penelitian dan pengembangan berjalan sesuai ruang lingkup yang ditentukan, sehingga hasil yang diperoleh dapat dipertanggungjawabkan dengan baik.
% --- Metodologi Pengerjaan TA ---
\section{Metodologi}

Untuk menyelesaikan permasalahan pada tugas akhir ini, digunakan metodologi \textit{Software Development Life Cycle} (SDLC) sebagai kerangka kerja dalam proses pengembangan sistem pelacakan ITB Ultra-Marathon. SDLC merupakan serangkaian tahapan sistematis yang digunakan untuk merancang, membangun, menguji, dan memelihara perangkat lunak agar sesuai dengan kebutuhan pengguna \cite{pressman2014}. Metodologi ini membantu memastikan bahwa setiap tahap pengembangan dilakukan secara terstruktur dan menghasilkan sistem yang fungsional serta memenuhi tujuan yang telah ditentukan.

Model SDLC yang digunakan dalam penyelesaian tugas akhir ini adalah model \textit{Waterfall}. Pemilihan metodologi ini didasarkan pada batasan ruang lingkup sistem yang telah ditetapkan serta kebutuhan tahapan pengembangan yang sistematis dan terstruktur. Pengembangan dilakukan secara linear dan sekuensial, di mana setiap tahap harus diselesaikan sebelum melanjutkan ke tahap berikutnya, sehingga memudahkan pengukuran progres dan evaluasi hasil secara mendalam. Adapun tahapan-tahapan yang dilakukan pada metode ini adalah sebagai berikut:

\begin{enumerate}
    \item {Analisis Kebutuhan}
    Tahap awal ini dilakukan untuk mengidentifikasi seluruh kebutuhan sistem, baik fungsional maupun non-fungsional, melalui proses pengumpulan data dan analisis terhadap penyelenggaraan ITB Ultra-Marathon pada tahun-tahun sebelumnya. Selain itu, dilakukan juga wawancara dengan pihak panitia penyelenggara tahun sebelumnya untuk memahami kebutuhan pengguna serta kendala yang terjadi pada sistem pelacakan yang telah digunakan.

    \item {Perancangan Sistem}

    Berdasarkan hasil analisis kebutuhan, dilakukan perancangan arsitektur sistem dan \textit{user interface}. Fokus pada tahap ini adalah merancang \textit{wireframe} dan \textit{high-fidelity design} yang intuitif, serta merancang skema komunikasi antara \textit{front-end} dan \textit{back-end}.

    \item {Implementasi}

    Tahap ini merupakan proses penerjemahan rancangan ke dalam baris kode program. Pengembangan difokuskan pada pembangunan sisi \textit{front-end} menggunakan kerangka kerja yang mendukung perangkat Android. Tahapan ini memastikan setiap elemen visual dan fitur pelacakan dapat beroperasi sesuai dengan spesifikasi yang telah ditetapkan.

    \item {\textit{Deployment}}

    Pada tahap ini, aplikasi yang telah dibangun didistribusikan ke dalam lingkungan pengguna. Hal ini mencakup proses instalasi dan konfigurasi sistem pada perangkat Android yang akan digunakan oleh pihak-pihak terkait dalam simulasi atau kegiatan ITB Ultra-Marathon.

    \item {Pengujian}

    Tahap pengujian dilakukan untuk mengevaluasi fungsionalitas sistem yang telah dikembangkan serta memastikan bahwa seluruh fitur dapat berjalan sesuai dengan kebutuhan yang telah ditentukan. Pengujian juga dilakukan untuk mengidentifikasi kesalahan pada sistem dan menilai kesesuaian antarmuka terhadap penggunaan aplikasi.

    \item {Evaluasi}

    Tahap akhir ini dilakukan untuk meninjau kembali apakah sistem yang telah dibangun telah memenuhi seluruh kriteria keberhasilan dan menjawab rumusan masalah. Evaluasi dilakukan dengan membandingkan hasil akhir terhadap target efisiensi manajemen perlombaan dan kualitas pengalaman pengguna saat mengoperasikan aplikasi.
\end{enumerate}

% --- Sistematika Penulisan ---
\section{Sistematika Penulisan}
Sistematika penulisan laporan tugas akhir ini disusun secara sistematis untuk memberikan gambaran menyeluruh mengenai alur penelitian dan pengembangan aplikasi. Adapun struktur penulisan laporan adalah sebagai berikut:

\begin{enumerate}
    \item Bab I Pendahuluan: Menguraikan latar belakang permasalahan, rumusan masalah, tujuan penelitian, batasan masalah, metodologi pengembangan, serta sistematika penulisan laporan.
    \item Bab II Studi Literatur: Memaparkan landasan teori dan kajian pustaka yang relevan, mencakup konsep \textit{ultra-marathon}, pengembangan \textit{front-end}, teknologi pelacakan, serta literatur mengenai sistem pendukung perlombaan lari jarak jauh.
    \item Bab III Analisis: Menjelaskan proses analisis masalah yang terjadi pada sistem pelacakan ITB Ultra-Marathon serta identifikasi kebutuhan fungsional dan non fungsional dari berbagai perspektif pengguna, yaitu peserta, panitia, tim support, dan penonton.
    \item Bab IV Perancangan: Menguraikan perancangan solusi yang diusulkan, meliputi arsitektur sistem secara menyeluruh, desain antarmuka pengguna, serta perancangan integrasi antara sisi \textit{front-end} dengan \textit{back-end}.
    \item Bab V Implementasi: Menjelaskan proses realisasi atau pengkodean antarmuka aplikasi berdasarkan rancangan yang telah dibuat, termasuk penggunaan teknologi, kerangka kerja (\textit{framework}), serta pustaka (\textit{library}) yang digunakan untuk memvisualisasikan data pelacakan.
    \item Bab VI Evaluasi: Berisi laporan mengenai metode pengujian sistem, persiapan skenario pengujian, hasil pengujian, serta evaluasi kinerja antarmuka dalam menyajikan data secara \textit{real-time}.
    \item Bab VII Kesimpulan dan Saran: Menyampaikan kesimpulan akhir dari seluruh proses pengembangan berdasarkan tujuan yang telah ditetapkan serta saran bagi pengembangan aplikasi di masa mendatang.
    \item Daftar Pustaka: Mencantumkan daftar referensi, buku, jurnal, dan sumber data lainnya yang digunakan sebagai acuan dalam penyusunan laporan ini.
    \item Lampiran: Menyertakan dokumen pendukung tambahan seperti kode program utama, dokumentasi pengujian, serta data pelengkap lainnya.
\end{enumerate}